% ============================================================
%  Basado en la plantilla de la Practica 6 (mismos colores/estilos)
% ============================================================
\documentclass[11pt,a4paper]{article}

% ---------- Idioma y codificación ----------
\usepackage[utf8]{inputenc}
\usepackage[T1]{fontenc}
\usepackage[spanish,es-noshorthands]{babel}

% ---------- Tipografía y layout ----------
\usepackage{lmodern}
\usepackage[sfdefault,light]{roboto}
\usepackage{geometry}
\geometry{a4paper, top=3.1cm, bottom=2.8cm, left=2.4cm, right=2.4cm, headheight=1.6cm}
\usepackage{ragged2e}
\usepackage{enumitem}
\usepackage{graphicx}
\usepackage{float}
\usepackage{booktabs}
\usepackage{array}
\usepackage[most]{tcolorbox}
\usepackage{tikz}
\usetikzlibrary{shadows,calc}
\usepackage{titlesec}
\usepackage{fancyhdr}
\usepackage{fontawesome5}
\usepackage[colorlinks=true,breaklinks=true]{hyperref}
\usepackage{caption}
\usepackage{microtype}
\usepackage{xurl}
\usepackage{listings}

% ---------- Paleta de colores (tema "coolnight" de tu terminal Kitty) ----------
\definecolor{cNavy}{HTML}{011423}
\definecolor{cTealDeep}{HTML}{064B5A}
\definecolor{cTeal}{HTML}{097289}
\definecolor{cBlue}{HTML}{0FC5ED}
\definecolor{cCyan}{HTML}{0FC5ED}
\definecolor{cMint}{HTML}{44FFB1}
\definecolor{cPurple}{HTML}{A277FF}
\definecolor{cRed}{HTML}{E52E2E}
\definecolor{cBg}{HTML}{EAF6F8}
\definecolor{cText}{HTML}{01111E}
\definecolor{cOutBg}{HTML}{F4F4F4}

\hypersetup{
  linkcolor=cTealDeep,
  citecolor=cTealDeep,
  urlcolor=cBlue,
  pdftitle={Práctica 7 -- Algoritmos de Clasificación en R},
  pdfauthor={Miguel Deza}
}

\color{cText}
\emergencystretch=2.5em

% ---------- Encabezado / pie de página ----------
\pagestyle{fancy}
\fancyhf{}
\renewcommand{\headrulewidth}{0pt}
\renewcommand{\footrulewidth}{0pt}
\fancyhead[L]{\small\color{cTealDeep}\textbf{\faGraduationCap\ Maestría en Ciencia de Datos}}
\fancyhead[R]{\small\color{cTeal}Práctica 7 -- Clasificación en R}
\fancyfoot[C]{%
  \begin{tikzpicture}[remember picture]
    \node[circle, fill=cTealDeep, minimum size=1.7em, inner sep=0pt] (pg) {\color{white}\small\thepage};
  \end{tikzpicture}%
}
\renewcommand{\headrule}{{\color{cCyan}\hrule height 1.4pt}\vspace{1pt}{\color{cBlue}\hrule height 0.5pt}}

% ---------- Estilo de secciones ----------
\titleformat{\section}
  {\normalfont\Large\bfseries\color{cTealDeep}}
  {\tikz[baseline=-0.6ex]{\node[fill=cCyan, text=white, rounded corners=1.5pt, inner sep=3pt, font=\bfseries]{\thesection};}}
  {0.6em}{}
  [{\vspace{-0.3em}\color{cCyan}\titlerule[1.2pt]\vspace{0.3em}}]

\titleformat{\subsection}
  {\normalfont\large\bfseries\color{cTeal}}
  {\color{cTeal}\thesubsection}{0.6em}{}

\titlespacing*{\section}{0pt}{1.4em}{0.8em}
\titlespacing*{\subsection}{0pt}{1.1em}{0.5em}

% ---------- Figura sin marco ----------
\newcommand{\figuraMarcada}[3]{%
  \begin{figure}[H]
    \centering
    \includegraphics[width=#2\linewidth]{#1}
    \caption{#3}
  \end{figure}
}

\captionsetup{labelfont={bf,color=cTealDeep}, font=small, textfont=it}

% ---------- Bloques de código R y salida de consola ----------
\lstdefinestyle{styleR}{
  language=R,
  backgroundcolor=\color{cBg},
  basicstyle=\ttfamily\small\color{cText},
  keywordstyle=\color{cTealDeep}\bfseries,
  commentstyle=\color{gray}\itshape,
  stringstyle=\color{cTeal},
  numbers=none,
  breaklines=true,
  frame=l,
  framerule=2.2pt,
  rulecolor=\color{cCyan},
  xleftmargin=8pt,
  framexleftmargin=6pt,
  aboveskip=6pt, belowskip=6pt,
  showstringspaces=false
}

\lstdefinestyle{styleOut}{
  backgroundcolor=\color{cOutBg},
  basicstyle=\ttfamily\scriptsize\color{cText},
  breaklines=true,
  frame=l,
  framerule=2.2pt,
  rulecolor=\color{cTealDeep},
  xleftmargin=8pt,
  framexleftmargin=6pt,
  aboveskip=2pt, belowskip=6pt,
  showstringspaces=false
}

\newcommand{\codigoR}[1]{\lstinputlisting[style=styleR]{#1}}

% ============================================================
\begin{document}

% -------------------- PORTADA --------------------
\begin{titlepage}
\newgeometry{margin=0cm}
\begin{tikzpicture}[remember picture, overlay]
  \fill[cNavy] (current page.south west) rectangle (current page.north east);

  \fill[cPurple, opacity=0.28] (current page.north east) ++(-1,1) circle (3.2);
  \fill[cCyan, opacity=0.22] (current page.north east) ++(-2.6,0.4) circle (2.1);
  \fill[cTealDeep, opacity=0.55] (current page.south west) ++(1,-1) circle (4.0);
  \fill[cMint, opacity=0.16] (current page.south west) ++(2.4,-0.2) circle (2.3);

  % ---- Icono insignia ----
  \node[circle, fill=cTealDeep, draw=cCyan, line width=1.3pt, minimum size=2.5cm] (logo) at ($(current page.north)+(0,-2.9)$) {};
  \node at (logo) {\color{cMint}\fontsize{28}{28}\selectfont\faIcon{r-project}};

  \node[anchor=north, text width=14cm, align=center] at ($(current page.north)+(0,-4.7)$) {
    {\color{cCyan}\bfseries\fontsize{12}{14}\selectfont \faDatabase\ \ MAESTRÍA EN CIENCIA DE DATOS}
  };

  \node[anchor=center, text width=15cm, align=center] at ($(current page.center)+(0,1.2)$) {
    {\color{white}\bfseries\fontsize{40}{44}\selectfont Práctica 7}\\[0.35em]
    {\color{cMint}\bfseries\fontsize{20}{24}\selectfont \faIcon{r-project}\ Algoritmos de Clasificación en R}
  };

  \draw[cCyan, line width=1.4pt] ($(current page.center)+(-3.6,-0.55)$) -- ++(7.2,0);

  \node[anchor=center, text width=13cm, align=center] at ($(current page.center)+(0,-1.3)$) {
    {\color{white!85!black}\fontsize{11}{14}\selectfont
     LDA, CART, kNN, SVM lineal y Random Forest sobre niveles de obesidad\\
     en personas de México, Perú y Colombia (Kaggle), comparados con \textbf{caret}}
  };

  \node[anchor=south, text width=14cm, align=center] at ($(current page.south)+(0,2.6)$) {
    \begin{tabular}{r l}
      \color{cCyan}\bfseries \faUser\ Elaborado por: & \color{white} Miguel Deza \\[3pt]
      \color{cCyan}\bfseries \faGraduationCap\ Programa: & \color{white} Maestría en Ciencia de Datos \\[3pt]
      \color{cCyan}\bfseries \faBook\ Curso: & \color{white} Introducción a Ciencia de Datos \\[3pt]
      \color{cCyan}\bfseries \faChalkboardTeacher\ Docente: & \color{white} Yessenia Deysi Yari Ramos \\[3pt]
      \color{cCyan}\bfseries \faIcon{calendar-alt}\ Fecha: & \color{white} 21 de septiembre de 2026 \\
    \end{tabular}
  };
\end{tikzpicture}
\end{titlepage}
\restoregeometry

% -------------------- DESARROLLO --------------------
\section{\faCogs\ \ Desarrollo}

\subsection{\faBoxOpen\ \ 1. Cargar las bibliotecas necesarias}
\begin{lstlisting}[style=styleR]
library(caret)
library(MASS)          # LDA
library(rpart)          # CART
library(rpart.plot)
library(e1071)           # SVM (y soporte interno de caret)
library(kernlab)          # motor de SVM lineal usado por caret
library(randomForest)
library(corrplot)
library(ggplot2)
\end{lstlisting}
Todo cargó bien.

\subsection{\faFileImport\ \ 2. Cargar el conjunto de datos (CSV)}
Busqué en \url{https://www.kaggle.com/datasets?tags=13302-Classification} un dataset de mi interés e inicié sesión en Kaggle para descargarlo directamente. Elegí \textbf{Obesity Levels} (\texttt{fatemehmehrparvar/obesity-levels}): estimación del nivel de obesidad de personas de \textbf{México, Perú y Colombia} según hábitos alimenticios y actividad física.
\begin{lstlisting}[style=styleR]
datos <- read.csv("obesidad.csv", stringsAsFactors = TRUE)
niveles_obesidad <- c("Insufficient_Weight", "Normal_Weight",
                       "Overweight_Level_I", "Overweight_Level_II",
                       "Obesity_Type_I", "Obesity_Type_II", "Obesity_Type_III")
datos$NObeyesdad <- factor(datos$NObeyesdad, levels = niveles_obesidad)
str(datos)
\end{lstlisting}
\begin{lstlisting}[style=styleOut]
'data.frame': 2111 obs. of  17 variables:
 $ Age                           : num  21 21 23 27 22 29 23 22 24 22 ...
 $ Gender                        : Factor w/ 2 levels "Female","Male"
 $ Height                        : num  1.62 1.52 1.8 1.8 1.78 ...
 $ Weight                        : num  64 56 77 87 89.8 53 55 53 64 68 ...
 $ CALC                          : Factor w/ 4 levels "Always","Frequently",...
 $ FAVC                          : Factor w/ 2 levels "no","yes"
 $ FCVC                          : num  2 3 2 3 2 2 3 2 3 2 ...
 $ NCP                           : num  3 3 3 3 1 3 3 3 3 3 ...
 $ SCC                           : Factor w/ 2 levels "no","yes"
 $ SMOKE                         : Factor w/ 2 levels "no","yes"
 $ CH2O                          : num  2 3 2 2 2 2 2 2 2 2 ...
 $ family_history_with_overweight: Factor w/ 2 levels "no","yes"
 $ FAF                           : num  0 3 2 2 0 0 1 3 1 1 ...
 $ TUE                           : num  1 0 1 0 0 0 0 0 1 1 ...
 $ CAEC                          : Factor w/ 4 levels "Always","Frequently",...
 $ MTRANS                        : Factor w/ 5 levels "Automobile","Bike",...
 $ NObeyesdad                    : Factor w/ 7 levels "Insufficient_Weight",...
\end{lstlisting}
2111 personas, 16 predictores (edad, género, altura, peso, hábitos de alimentación/actividad física y transporte) y la clase \texttt{NObeyesdad} con 7 niveles de obesidad.

\subsection{\faSearch\ \ 3. Análisis exploratorio de datos}
\begin{lstlisting}[style=styleR]
summary(datos)
table(datos$NObeyesdad)
colSums(is.na(datos))
cor(datos[, c("Age","Height","Weight","FCVC","NCP","CH2O","FAF","TUE")])
\end{lstlisting}
\begin{lstlisting}[style=styleOut]
Balance de clases:
Insufficient_Weight       Normal_Weight  Overweight_Level_I Overweight_Level_II
                272                 287                 290                 290
     Obesity_Type_I     Obesity_Type_II    Obesity_Type_III
                351                 297                 324

Proporcion (%): 12.9  13.6  13.7  13.7  16.6  14.1  15.3

Valores faltantes: 0 en las 17 columnas
\end{lstlisting}
Clase \textbf{bien balanceada} (12.9\%--16.6\% por nivel) y \textbf{sin valores faltantes}. En la correlación entre numéricas, \texttt{Height}--\texttt{Weight} es la más fuerte (0.46), seguida de \texttt{Height}--\texttt{FAF} (0.29) y \texttt{Age}--\texttt{TUE}, la única negativa relevante (-0.30: las personas mayores usan menos dispositivos electrónicos).

\subsection{\faChartBar\ \ 4. Visualizar el conjunto de datos}
\figuraMarcada{img/grafico-balance-clases.png}{0.72}{Balance de clases. Uso una paleta secuencial (no categórica) porque \texttt{NObeyesdad} es ordinal: a mayor nivel, más obesidad.}

\figuraMarcada{img/grafico-boxplots-variables.png}{0.95}{Peso por nivel de obesidad: la relación es perfectamente monótona, tal como debe ser por construcción de las categorías.}

\figuraMarcada{img/grafico-correlacion.png}{0.62}{Correlación entre los 8 predictores numéricos.}

\subsection{\faTools\ \ 5. Preprocesar los datos}
\begin{lstlisting}[style=styleR]
set.seed(7)
idx_train <- createDataPartition(datos$NObeyesdad, p = 0.80, list = FALSE)
train <- datos[idx_train, ]
test  <- datos[-idx_train, ]
\end{lstlisting}
\begin{lstlisting}[style=styleOut]
Filas train / test:
train  test
 1691   420
\end{lstlisting}
No hizo falta imputar nada (sin NA). Las 8 columnas categóricas (\texttt{Gender}, \texttt{CALC}, \texttt{FAVC}, \texttt{SCC}, \texttt{SMOKE}, \texttt{family\_history\_with\_overweight}, \texttt{CAEC}, \texttt{MTRANS}) ya quedaron como \texttt{factor}; \texttt{train()} las convierte a variables \textit{dummy} automáticamente vía la fórmula. Centré y escalé las numéricas dentro de cada \texttt{train()} (\texttt{preProcess = c("center","scale")}) porque kNN y SVM son sensibles a la escala.

\subsection{\faPlay\ \ 6. Ejecutar los algoritmos}
\begin{lstlisting}[style=styleR]
control <- trainControl(method = "repeatedcv", number = 10, repeats = 3)
metric  <- "Accuracy"
prep    <- c("center", "scale")

fit.lda  <- train(NObeyesdad ~ ., data = train, method = "lda",
                   metric = metric, trControl = control, preProcess = prep)
fit.cart <- train(NObeyesdad ~ ., data = train, method = "rpart",
                   metric = metric, trControl = control, tuneLength = 8)
fit.knn  <- train(NObeyesdad ~ ., data = train, method = "knn",
                   metric = metric, trControl = control, preProcess = prep, tuneLength = 10)
fit.svm  <- train(NObeyesdad ~ ., data = train, method = "svmLinear",
                   metric = metric, trControl = control, preProcess = prep)
fit.rf   <- train(NObeyesdad ~ ., data = train, method = "rf",
                   metric = metric, trControl = control)
\end{lstlisting}
Los cinco modelos usan la misma validación cruzada: \textbf{10 pliegues $\times$ 3 repeticiones}. Al entrenar LDA, R avisó \texttt{"variables are collinear"} -- con 8 columnas categóricas convertidas a \textit{dummies} (algunas con niveles muy raros, como \texttt{CALC="Always"} con solo 1 caso), el sistema de ecuaciones de LDA queda casi singular; no impide el ajuste, pero es una limitación real de LDA frente a datos con muchas variables \textit{dummy} poco frecuentes.

\subsection{\faChartBar\ \ 7. Evaluar los algoritmos}
\begin{lstlisting}[style=styleR]
resultados <- resamples(list(LDA=fit.lda, CART=fit.cart, kNN=fit.knn,
                              SVM=fit.svm, RF=fit.rf))
summary(resultados)
\end{lstlisting}
\begin{lstlisting}[style=styleOut]
Accuracy (media de 30 evaluaciones)
LDA  0.8965   CART 0.7765   kNN 0.7905   SVM 0.9515   RF 0.9645
Kappa (media)
LDA  0.8791   CART 0.7392   kNN 0.7553   SVM 0.9433   RF 0.9585
\end{lstlisting}
\figuraMarcada{img/grafico-comparacion-cv.png}{0.85}{Accuracy y Kappa de las 30 evaluaciones de validación cruzada, por algoritmo. Jerarquía clara: RF $>$ SVM $>$ LDA $>$ kNN $>$ CART.}

\begin{lstlisting}[style=styleR]
pred.lda <- predict(fit.lda, test); pred.rf <- predict(fit.rf, test)
confusionMatrix(pred.lda, test$NObeyesdad)
# ... (una matriz de confusion 7x7 por modelo)
\end{lstlisting}

\begin{table}[H]
\centering
\small
\begin{tabular}{lrrrrr}
\toprule
\textbf{Modelo} & \textbf{Accuracy CV} & \textbf{Accuracy test} & \textbf{Kappa test} & \textbf{Sensib. macro} & \textbf{Especif. macro} \\
\midrule
LDA            & 0.896 & 0.879          & 0.858          & 0.875 & 0.980 \\
CART           & 0.777 & 0.783          & 0.747          & 0.789 & 0.964 \\
kNN            & 0.790 & 0.819          & 0.789          & 0.815 & 0.970 \\
SVM lineal     & 0.951 & 0.943          & 0.933          & 0.941 & 0.991 \\
Random Forest  & \textbf{0.965} & \textbf{0.969} & \textbf{0.964} & \textbf{0.969} & \textbf{0.995} \\
\bottomrule
\end{tabular}
\caption{Tabla comparativa final sobre el conjunto de prueba (420 personas). Sensibilidad/especificidad son el promedio macro de las 7 clases.}
\end{table}

\textbf{Random Forest es claramente el mejor} (96.9\% de accuracy en prueba, Kappa 0.964), seguido de SVM lineal (94.3\%). CART, el modelo más simple, es el más débil (78.3\%) -- señal de que las fronteras entre niveles de obesidad no son un solo corte rectangular, sino que se benefician de un ensamble de árboles o un margen lineal de alta dimensión.

\figuraMarcada{img/grafico-arbol-cart.png}{0.98}{Árbol CART final: 6 niveles de profundidad, dividiendo primero por \texttt{Weight} y luego por \texttt{Height} y \texttt{Gender}.}

\figuraMarcada{img/grafico-importancia-rf.png}{0.85}{Importancia de variables (Random Forest, \textit{Mean Decrease Gini}): \texttt{Weight} y \texttt{Height} dominan por un margen amplio.}

Tanto el árbol CART como Random Forest coinciden en que \texttt{Weight} y \texttt{Height} son los predictores más fuertes -- razonable, dado que la clase \texttt{NObeyesdad} está construida a partir del IMC (peso/altura$^2$).

\subsection{\faMagic\ \ 8. Hacer predicciones}
Con Random Forest (mejor modelo), sobre una muestra del set de prueba:
\begin{lstlisting}[style=styleR]
muestra <- test[sample(nrow(test), 6), ]
predict(fit.rf, muestra)
\end{lstlisting}
\begin{lstlisting}[style=styleOut]
                 Real          Prediccion
1      Obesity_Type_I      Obesity_Type_I
2    Obesity_Type_III    Obesity_Type_III
3 Insufficient_Weight Insufficient_Weight
4 Overweight_Level_II Overweight_Level_II
5      Obesity_Type_I      Obesity_Type_I
6 Overweight_Level_II Overweight_Level_II
\end{lstlisting}
6 de 6 aciertos.

Y sobre una persona nueva (hipotética, perfil peruano típico):
\begin{lstlisting}[style=styleR]
persona_nueva <- data.frame(
  Age = 27, Gender = "Male", Height = 1.72, Weight = 88,
  CALC = "Sometimes", FAVC = "yes", FCVC = 2, NCP = 3, SCC = "no", SMOKE = "no",
  CH2O = 2, family_history_with_overweight = "yes", FAF = 1, TUE = 2,
  CAEC = "Sometimes", MTRANS = "Public_Transportation")
predict(fit.rf, persona_nueva)
predict(fit.rf, persona_nueva, type = "prob")
\end{lstlisting}
\begin{lstlisting}[style=styleOut]
[1] Overweight_Level_II
  Insufficient_Weight Normal_Weight Overweight_Level_I Overweight_Level_II
1                   0         0.002               0.07               0.854
  Obesity_Type_I Obesity_Type_II Obesity_Type_III
1          0.074               0                0
\end{lstlisting}
Con 1.72\,m y 88\,kg (IMC $\approx$ 29.7), el modelo predice \texttt{Overweight\_Level\_II} con 85.4\% de probabilidad -- consistente con el rango de IMC de esa categoría.

\end{document}
